% Ad Quintum Fratrem 1.1 (ad-quintum-fratrem-01-01) --- English translation
% Translated by Claude (Anthropic) from the Latin (Perseus).
% Style guide: STYLE.md.

\ciceroLetterOpener{Marcus to his brother Quintus, greetings.}

\ciceroSection{1} Although I had no doubt that this letter would be outstripped, by the speed of rumour itself, by many messengers, and that you would hear earlier from others that a third year has been added to my longing and to your labour, yet I thought it fitting that the news of this trouble be carried to you by me also. For in earlier letters --- not in one but in many --- when the matter was already despaired of by others, I yet brought you hope of an early withdrawal, not only that I might delight you as long as possible with a pleasing opinion, but also because so great an effort was being applied both by us and by the praetors that I did not despair of the thing being brought off.

\ciceroSection{2} Now, since it has so happened that neither could the praetors with their resources nor we with our zeal accomplish anything, it is altogether hard not to feel it heavily; but yet it is not fitting that our spirits, exercised in the greatest enterprises both of doing and of enduring, be broken and weakened by trouble. And since men ought to bear most heavily what has been brought on by their own fault, there is something in this matter that I must bear more heavily than you. For it has come about by my fault, contrary to what you had urged me as you were setting out and by letter, that no successor was sent in the previous year. Which I, while consulting the safety of the allies, while resisting the impudence of certain businessmen, while seeking to increase our glory by your virtue, did unwisely --- especially since I committed the matter to such a turn that the second year of yours could entail a third.

\ciceroSection{3} Since I confess that this is my error, it is the part of your wisdom and humanity to take care, and to bring it about, that what was less wisely foreseen by me be corrected by your diligence. And if you stir yourself up vehemently to all parts of having a good name --- so as no longer to compete with others but with yourself --- if you rouse all your mind, your care, your thought to the desire of an excellent praise in everything, believe me, this one year added to your labour will bring me, and our descendants too, the joy of many years.

\ciceroSection{4} Wherefore I ask you first this: do not draw in and let your spirit fall, do not allow yourself to be overwhelmed, as by a tide, by the size of the business; but rather rouse yourself and stand against it, or rather go forward to meet the business of your own accord. For you are not handling that part of the commonwealth in which fortune is master, but one in which reason and diligence can do most. If I saw your command extended while you were managing some great and dangerous war, I should tremble in my mind, since I should understand that fortune's power, too, had at the same time been extended over us.

\ciceroSection{5} As things are, however, that part of the commonwealth has been entrusted to you in which fortune holds either no part or a very small one, and which seems to me to depend wholly on your virtue and moderation of mind. No plots of the enemy, I take it, no contest of battle, no defection of allies, no want of pay or of grain, no mutiny of the army, do we fear --- the kind of things that have very often befallen the wisest men, so that, as the best helmsmen could not overcome the force of a storm, they could not overcome the impulse of fortune. To you has been given the highest peace, the highest tranquillity --- but in such a way that even peace can sink a sleeping helmsman, and yet delight a wakeful one.

\ciceroSection{6} For that province consists, first, of that kind of allies which is, of every kind of men, the most humane; second, of that kind of citizens who, either because they are publicans, are bound to me by the highest connection; or because they so do business as to be wealthy, think that they have their fortunes safe by the kindness of our consulship.

\ciceroSection{7} ``But,'' you will say, ``among these very men there arise heavy controversies; many injuries are born; great contests follow.'' As if I were thinking that you do not have any small business on your hands! I understand that the business is very great and of the highest counsel; but remember, I think this business is more a matter of counsel than of fortune. For what is it to keep within bounds those whom you preside over, if you keep yourself within bounds? Let that be a great and difficult thing for others, hard as it is. For you it has always been very easy --- and ought indeed to be so, whose nature is such that, even without learning, it would seem to have been able to be moderate; and the learning has been brought to bear which could even cultivate the most vicious nature. You, when you withstand money, when you withstand pleasure, when you withstand the desire of every kind of thing (as you do), it will be a danger, I take it, that you cannot keep down a wicked businessman or a slightly more grasping publican. As for the Greeks, they will so look at you, living thus, as to think that some man out of the memory of the annals or even some divine man from heaven has come down into the province.

\ciceroSection{8} I write these things now not in order that you should do them, but that you may rejoice that you do them and have done them. For it is a fine thing that the highest command for three years in Asia has been such that no statue, no painting, no vase, no garment, no slave, no person's beauty, no terms of money --- things in which that province abounds --- has drawn you from the highest integrity and self-control.

\ciceroSection{9} What can be found so excellent or so to be wished as that this virtue, this moderation of mind, this temperance, should not lie hidden in darkness, nor be put away, but be set in the light of Asia, in the eyes of the most distinguished province, in the ears of all peoples and nations? Men are not terrified by your journeys, nor exhausted by expense, nor stirred up by your arrival; the highest gladness, both public and private, comes wherever you have come, since the city seems to have received a guardian, not a tyrant, and the house a guest, not a plunderer.

\ciceroSection{10} In these matters, however, experience itself has now surely instructed you that it is by no means enough that you yourself have these virtues; you must also look about diligently, so that, in this guardianship of the province, you may be seen to vouch to the allies and to the citizens and to the commonwealth not for yourself alone but for all the agents of your command. Although you have legates of the kind who, on their own account, will have a regard for their own dignity. Of these, in honour and in dignity and in age, Tubero stands first; whom (especially since he is writing history) I think can choose, from his own annals, many men whom he wishes and is able to imitate. Allienus is ours both in spirit and in goodwill, and indeed in his way of life by imitation. As for Gratidius --- what shall I say? --- whom I know with certainty so labours about his own reputation that, on account of his brotherly love for us, he labours also about ours.

\ciceroSection{11} A quaestor you have, not chosen by your judgment but the man whom the lot gave. He must both be moderate of his own accord and obey your principles and instructions. If any of these should perchance be more sordid, you would bear with him only so long as he, on his own account, neglected those laws by which he was bound; not that the power which you have entrusted to him for dignity, he should turn to gain. For it does not at all please me, especially since these morals have already inclined too much to excessive softness and to favour-seeking, that you should pry into all squalid matters and shake out every man one by one; rather, that according to each man's faith, so much should be entrusted to each. And among these, those whom the commonwealth itself has given you as companions and helpers in public business, you will keep them within the bounds I have just prescribed.

\ciceroSection{12} As to those, however, whom you have wished to be with you out of household intimacy or out of necessary attendances, who are wont to be called as if from the praetor's cohort --- of these we are answerable not only for all the deeds but also for all the words. But you have those with you whom you can easily love when they act rightly, and most easily restrain when they consult your reputation too little. By them, when you were inexperienced, your generosity might have been deceived; for, as each man is the best, so most reluctantly does he suspect that others are wicked. Now indeed let this third year hold the same integrity as the previous, and an even more cautious and diligent one.

\ciceroSection{13} Let your ears be such as are thought to hear what they hear, not such as one whispers into feignedly and pretendedly for gain. Let your seal-ring be not as some implement, but as you yourself are: not the agent of another's will, but the witness of your own. Let the herald be of that number which our ancestors wished it to be: men whom they appointed to that role not as a kindness but as a labour and duty, and not without thought, except to their own freedmen, whom they commanded little differently from slaves. Let the lictor be the agent not of his own mildness but of yours; and let the fasces and axes carry before greater badges of dignity than of power. In short, let it be known to the whole province that to you the safety, the children, the reputation, the fortunes of all whom you preside over are most dear. Finally, let this be the opinion: that you will be an enemy not only to those who have taken anything but to those also who have given, if you have learned of it. Nor indeed will anyone give, when this is clear: that nothing is wont to be obtained from you through those who pretend to have great influence with you.

\ciceroSection{14} And yet my speech is not of such a kind that I want you to be either too hard with your own people or too suspicious. For if there is anyone of them who in two years has never come into the suspicion of greed --- as I both hear and (because I know them) judge of Caesius and Chaerippus and Labeo --- there is nothing that I do not think can most rightly both be entrusted and credited to them, and to any other man of the same kind. But if there is anyone in whom you have already taken offence, of whom you have suspected something, to him you would credit nothing, and entrust no part of your reputation.

\ciceroSection{15} In the province itself, however, if you have found someone who has come deeply into your familiarity, who was previously unknown to us --- consider how far he is to be trusted. Not but that there can be many honest men in the provinces; but it is permitted to hope this; it is dangerous to judge it. For each man's nature is covered by many veils of pretence and is, as it were, drawn over by certain curtains. The forehead, the eyes, the face very often lie; speech, indeed, lies most often. Wherefore how can you find from that kind of men those who, drawn by desire of money, lack all those things from which we cannot be torn away, and who love you, a stranger, from the heart and not pretend it for their own gain? It seems to me very great --- especially if the same men love almost no private man, but always love every praetor. From this kind, if you have happened to know any man more loving of you than of the moment (for it has been possible), gladly enrol him in your number. But if you do not perceive that, no kind of friendship is more to be guarded against, since they know all the paths of money and do everything for money's sake, and do not care to consult the reputation of one with whom they are not going to live.

\ciceroSection{16} Even of the Greeks themselves certain familiarities are to be diligently guarded against, except those of a very few men --- if there are any worthy of old Greece. Most of them, alas, are deceitful and light, and trained by long servitude to excessive flattery. Whom, all of them, I say should be welcomed liberally; the best of them should be joined by hospitality and friendship; their excessive familiarities are neither so faithful (for they do not dare to oppose our wishes), and they envy not only ours but their own.

\ciceroSection{17} Now if I would be cautious and diligent in such matters --- in which I fear I am even too hard --- with what spirit do you suppose I am toward slaves? Whom indeed we ought to govern in every place, but especially in the provinces. Of which kind much could be prescribed, but this is the briefest and most easily kept: that they so conduct themselves on those Asian journeys as if you were going by the Appian Way, and that they think it makes no difference whether they have come to Tralles or to Formiae. And if any of the slaves is outstandingly faithful, let him be in domestic and private matters; in matters that pertain to the duty of your command and to any part of the commonwealth, let him touch nothing. For many things which can rightly be entrusted to faithful slaves nevertheless, for the sake of avoiding talk and reproach, must not be entrusted.

\ciceroSection{18} But somehow my speech has slipped into the manner of a manual of instruction, although that was not my purpose at the start. For what could I prescribe to a man who, especially in this kind, is, as I see, not inferior to me in prudence and even superior in experience? But yet I thought that, if my authority should be added to what you do, it would be the more pleasing to you yourself. Wherefore let these be the foundations of your dignity: first, your own integrity and self-control; next, the modesty of all who are with you, a careful and diligent selection in your familiarities both with provincials and with Greeks, a grave and steady discipline of the household.

\ciceroSection{19} Things which, while they are honourable in our private and daily affairs, in so great a command, with morals so depraved, with a province so corrupting, must necessarily seem divine. This training and this discipline can sustain in things that have to be set down and decreed that severity which you have used in matters from which we have undertaken some quarrels (with great joy on my part). Unless you suppose that I am moved by the complaints of some Paconius (a man not even Greek but a Mysian or rather a Phrygian) or by the voice of Tuscenius (a frantic and sordid man), out of whose foulest jaws you snatched a most dishonourable greed with the highest equity.

\ciceroSection{20} These and other rulings of yours full of severity in that province we should not easily sustain without the highest integrity. Wherefore let there be the highest severity in pronouncing the law, provided that it not be varied by favour but be kept consistent. But it is of small account that the law be pronounced consistently and diligently by you yourself, unless the same is done by those to whom you have ceded any part of that office. To me it seems that there is not much variety of business in administering Asia, and that the whole is sustained chiefly by the saying of justice. In which (especially of provincial knowledge) the principle itself lies open; what must be applied is steadfastness and gravity which can withstand not only favour but even suspicion.

\ciceroSection{21} There must also be added an easy bearing in hearing, mildness in deciding, diligence in giving satisfaction and in argument. By these things lately Gaius Octavius was most pleasant; in his court the first lictor was quiet, the herald was silent, every man spoke as often as he wished and as long as he wished. By all this perhaps he would have seemed too mild, were it not that this mildness defended that severity. The Sullan men were compelled to give back what they had taken away by force and fear. The men who in their magistracies had decreed wrongfully had to obey the same law themselves as private men. This severity of his would have seemed bitter, had it not been seasoned with many seasonings of humanity.

\ciceroSection{22} If this mildness is welcome at Rome, where there is so great arrogance, so unmeasured liberty, so endless licence of men, finally so many magistrates, so many helps, so great authority of the people and of the senate, how much more welcome may a praetor's courtesy be in Asia, in which so great a multitude of citizens, so great a multitude of allies, so many cities, so many states look to the nod of one man --- where there is no help, no place to complain, no senate, no public meeting? Wherefore it is the part of a great man, both naturally moderate and indeed schooled by learning and the best studies, so to bear himself in such great power that no other power may be missed by those over whom he presides.

\ciceroSection{23} The Cyrus of Xenophon was written not for the truth of history but for the image of just rule, in whom the highest gravity is joined by that philosopher with singular courtesy. Which books, indeed, our Africanus was not without reason wont not to lay out of his hands. For no duty of a diligent and moderate command has been passed over in them. And if those things were so cultivated by him who was never going to be a private citizen, in what manner must they be observed by those to whom command has been so given that they must give it back, and given by laws to which they must return?

\ciceroSection{24} And to me it seems that everything must be referred to this end by those who govern others: that those who shall be in their command be as happy as possible. That this is the most ancient principle with you, and was so from the beginning, when you first touched Asia, has been celebrated by constant rumour and by every man's report. It is the part not only of one who governs allies and citizens, but of one who governs slaves, and even mute herds, to serve the interests and the advantage of those whom he governs.

\ciceroSection{25} It is agreed by all, of this kind, that the highest diligence has been applied by you. No new debt has been contracted by the cities; many have been freed by you from the great and heavy old debt; cities ruined and almost deserted, including the most noble of Ionia, Samos, and one of Caria, Halicarnassus, have been set up again through your work. There are no seditions in the towns, no discords; you have provided that the cities be administered by the counsels of the best men. The brigandage of Mysia has been put down; massacre has been suppressed in many places; peace has been established throughout the province. And not only the brigandage of the roads and the fields, but also and much more the larger brigandage of towns and shrines, has been driven away. From the reputation, the fortunes, and the leisure of the wealthy, that bitterest handmaid of praetorial greed --- the false accusation --- has been removed. The expenses and tributes of the cities are equally borne by all who dwell in those cities' territories. Approach to you is most easy; your ears lie open to every man's complaint; no man's poverty and isolation has been shut out, not only from that public approach and the tribunal but not even from your house and chamber. In the whole of your command nothing is bitter, nothing cruel, and everything is full of clemency, mildness, humanity.

\ciceroSection{26} How great a benefit is yours --- that of an unfair and heavy aedilician tax, with great quarrels of mine, you have freed Asia! For if one nobleman complains openly of you, that by edicting that money should not be decreed for the games you have snatched HS 200,000 from him, what indeed would have been paid out, if money had been put forth in the name of all who hold games at Rome, as the practice now was? Yet we have suppressed these complaints of our men by this counsel: that --- in Asia, somehow, but at Rome praised with no slight wonder --- when the cities had decreed money for our temple and monument, and had done it both for our great services and for your supreme benefits, and that with the highest goodwill of theirs (and the law nominally provided that one might receive money for a temple and monument, and that what was given would not perish but be among the temple ornaments, so that it would seem given to the Roman people and the immortal gods rather than to me) --- yet that, in which lay dignity, the law, the doer's wish, I thought I should not accept, both for other reasons and that they --- to whom it was neither owed nor permitted --- should bear it with a more equable mind.

\ciceroSection{27} Wherefore lean with all your spirit and zeal on that course you have used so far: love those whom the senate and Roman people have committed and entrusted to your good faith and power; protect them by every reckoning; wish them as happy as possible. If the lot had set you over the Africans or the Spaniards or the Gauls --- savage and barbarous nations --- it would still be the part of your humanity to consider their advantage, and to serve their utility and safety. But when we are over a kind of men in whom is humanity itself, indeed from whom humanity is reckoned to have come to others, certainly to them of all men we ought to give it back, from whom we received it.

\ciceroSection{28} For I shall not now be ashamed to say this, especially in such a life and in such deeds in which no suspicion of sloth or levity can settle: that we have attained the things we have attained by those studies and arts which were handed down to us in the monuments and disciplines of Greece. Wherefore, beyond the common faith owed to all, we are seen to owe to that kind of men in particular this: that, with those whose teachings have schooled us, we should wish to display what we have learned from them.

\ciceroSection{29} And that prince of talent and learning, Plato, judged that commonwealths would be happy only when either learned and wise men should begin to govern them, or those who governed them should set their whole zeal on learning and wisdom; this joining of power and wisdom he reckoned could be a salvation to states. This perchance fell at some time to our whole commonwealth; now indeed surely it has fallen to that province of yours, that the man who in the receiving of learning, of virtue, of humanity, has had the most zeal and the most time from boyhood, holds the highest power in it.

\ciceroSection{30} Wherefore see that this year which has been added to your labour seem to have been extended at the same time for the safety of Asia. Since, in keeping you, Asia has been more fortunate than we in trying to bring you back, bring it about that the joy of the province may soften our longing. For if, in deserving the highest honours that I do not know any man has had, you were the most diligent of all, you ought to apply much greater diligence in keeping these honours.

\ciceroSection{31} Of that kind of honours, what I think I have written to you before. I have always reckoned them, if vulgar, cheap; if instituted for the moment, light. But if (as in fact has happened) they were paid to your merits, I judged that you should put much labour into keeping them. Wherefore, since you live in those cities with the highest command and power --- in which you see your virtues consecrated and placed among the gods --- in everything which you set up, decree, and do, you will think what you owe to such opinions of men, such judgments about you, such honours. And that will be of such a kind as to consult all men, to apply remedy to the troubles of men, to provide for safety, to wish to be called and held the parent of Asia.

\ciceroSection{32} And yet, to this wish of yours and to your diligence, the publicans bring great difficulty. If we oppose them, we shall cut off, both from us and from the commonwealth, an order which has best deserved of us and is bound through us to the commonwealth. But if we yield to them in everything, we shall let perish from the foundations those whose interest, as well as safety, we ought to consult. This is one --- if we wish to think honestly --- difficulty in your whole command. For to be self-restrained, to keep down all desires, to coerce one's own people, to keep an even reckoning of the law, to make oneself easily available in hearing causes and in admitting men, is a fine thing rather than a hard one; for it is set not in any labour but in a certain inclination and will of the mind.

\ciceroSection{33} What bitterness the matter of the publicans brings to the allies, we have learned from our own citizens, who were lately complaining (in the abolition of the tolls of Italy) not so much about the toll as about certain wrongs of the toll-collectors. Wherefore, since I have heard the complaints of citizens in Italy, I am not unaware of what may be happening to allies in the ends of the earth. To handle yourself there in such a way that you both satisfy the publicans (especially when their contracts have been badly let) and do not let the allies perish, seems the work of a divine kind of virtue --- which is yours. And first, what is most bitter to the Greeks --- that they are tributary --- ought not to seem so bitter, since without the command of the Roman people they have, by their own institutions, been so of themselves. They cannot recoil from the name ``publican,'' since they themselves could not have paid their tribute without a publican (which Sulla had distributed equally among them). That the Greeks are no milder in exacting tributes than our publicans can be understood from this: that the Caunians and all those of the islands which had been assigned by Sulla to the Rhodians lately fled to the senate that they might rather pay their tribute to us than to the Rhodians. Wherefore the name ``publican'' should be neither shuddered at by those who have always been tributary, nor abhorred by those who could not pay tribute by themselves, nor refused by those who have demanded it.

\ciceroSection{34} At the same time let Asia consider this also: that no calamity, either of foreign war or of domestic discords, would have been absent from her, if she were not held by this command. And since that command can in no way be retained without tributes, let her with an even mind redeem from herself a perpetual peace and quiet by some part of her own gains.

\ciceroSection{35} If they will then bear the kind and the name of publican without an unfair mind, the rest can seem milder to them by your counsel and prudence. They can, in making their pacts, look not at the censorial law but rather at the convenience of finishing the business and at the riddance of trouble. You yourself can do (as you have outstandingly done and are doing) this: bring to mind how great is the dignity of the publicans, how much we owe to that order, so that --- with command and force of authority and of the fasces removed --- you may join the publicans with the Greeks in favour and authority. And from those whom you have most served and who owe you everything, ask this: that by their easiness they allow us to keep and preserve that connection which we have with the publicans.

\ciceroSection{36} But why am I exhorting you in these matters, which you can do not only of your own accord without any man's prescriptions, but have already in great part accomplished? For the most honourable and greatest of the corporations do not stop giving us thanks daily; which to me is the more pleasant because the Greeks too do the same. To unite by goodwill those things which are diverse in interest, in utility, and almost in nature, is hard. But what is written above I have not written to instruct you (your prudence does not need any man's prescriptions), but the recollection of your virtue delighted me in the writing --- although in this letter I have been longer than I either wished or thought I would be.

\ciceroSection{37} There is one thing in which I shall not stop prescribing to you, nor allow you, so far as is in me, to be praised with an exception. For all who come from those parts so make mention of your virtue, integrity, humanity that, in the highest praises of you, they except one thing: irascibility. Which fault, although in this private and daily life it seems to be of a light and weak mind, yet is most ugly when bitterness of nature is added to the highest command. Wherefore I shall not undertake to set out for you now what is wont to be said about anger by the most learned men, since I do not wish to be too long, and you can easily learn it from the writings of many. That, however, which is proper to a letter --- that the man to whom it is written be informed about matters of which he is unaware --- I do not think should be passed over.

\ciceroSection{38} Almost all men report this to us: that, when anger is absent, you cannot be made more pleasant; but that, when the wickedness and perversity of any man have stirred you up, you are so inflamed in mind that everybody misses your usual humanity. Wherefore, since not so much some desire of glory as the very thing itself and fortune have led us into a way of life in which there will be perpetual talk of men about us, let us guard --- so far as we can effect and follow --- that no signal vice be said to have been in us. Nor do I now contend for what is perhaps difficult both in every nature and in our age now, to change one's mind and to pluck out suddenly anything that is deep-set in our manners. But I admonish you of this: that, if you cannot wholly avoid this --- since the mind is taken up by anger before reason can foresee that it might be taken up --- prepare yourself in advance, and meditate every day that anger must be resisted. And when it most stirs the mind, then you must most diligently restrain your tongue. Which virtue sometimes seems to me no less than not to be angry at all. For the latter is a sign not only of gravity but sometimes of slowness; but to moderate both mind and speech when you are angry, or even to be silent and to keep the mind's motion and grief in your own power --- although it is not the work of perfect wisdom, yet is the work of no slight talent.

\ciceroSection{39} And in this kind men report you have already become much more pleasant and milder. No more vehement excitements of mind of yours, no curses, no insults, are reported to us --- which both shrink from learning and humanity, and are above all contrary to command and dignity. For if angers are unappeasable, the bitterness is the highest; but if they can be entreated, the levity is the highest --- which yet, as among bad things, is to be preferred to bitterness.

\ciceroSection{40} Since the first year had the most talk on this rebuke (because, I suppose, men's wrongs, greed, and insolence happened to you beyond opinion and seemed unbearable), and the second year was much milder (because both habit and reason and, as I think, my letters too made you more patient and milder), the third year ought to be so corrected that no man can find fault with even the smallest thing.

\ciceroSection{41} And here at this point I deal with you, no longer with exhortation or prescriptions, but with brotherly entreaty: lay all your spirit, care, and thought into gathering praise from every quarter. If we stood at a moderate level of report and acclaim, no extraordinary thing would be required of you, nothing beyond the customary practice of others. But now, on account of the splendour and greatness of those things in which we have been busy, unless we attain the highest praise from that province, we scarcely seem able to escape the highest reproach. Such is our position: that all good men both favour us and demand and expect from us every diligence and virtue; and all the wicked --- since we have undertaken eternal war with them --- seem ready to be content with the slightest thing for finding fault.

\ciceroSection{42} Wherefore, since such a theatre as that of all Asia has been given to your virtues --- a theatre filled with the largest crowd, of the greatest size, of the most learned judgment, and so resounding by nature that all the way to Rome the demonstrations and the voices come back --- strive, I beg, and labour, not only to seem to have been worthy of these matters, but to have surpassed all of them by your skills.

\ciceroSection{43} And since chance has given me the urban administration of the commonwealth in office, and you the provincial: if my part yields to none, see that yours surpasses all the rest. Consider also: we are no longer labouring for some remaining and hoped-for glory, but fighting for what we have won --- which it was not so much our part to seek as to defend. If anything were separate from you for me, I would desire nothing more than the standing which has now been won for me. But now things stand thus: unless all your acts and words there answer our affairs here, I shall reckon that, by all my labours and by all my dangers (in all of which you were a sharer) I have attained nothing. If you above others helped that we might attain a most distinguished name, you will certainly above others labour that we may keep it. You must not use only the opinions and judgments of those who are now men, but those also who shall be in the future --- although those will be the truer judgment, freed from detraction and malice.

\ciceroSection{44} Lastly, you ought to consider this also: you do not seek glory for yourself alone. If that were so, you yet would not neglect it, especially since you have wished to consecrate the memory of your name with the most distinguished monuments. But it must be shared with me, and handed down to our children. In which it must be guarded against, if you should be too negligent, that you should seem to have consulted yourself little but even to have envied your own.

\ciceroSection{45} These things are not said in such a way that my speech may seem to have roused you sleeping, but rather to have urged you running. For you will perpetually do what you have done, so that all may praise your fairness, your temperance, your severity, your integrity. But me, on account of my singular love, an infinite desire of your glory holds. Although I judge this: since now Asia ought to be as well known to you as his own house to each man --- since to your highest prudence so great experience has been added --- there is nothing that pertains to praise which you do not most fully see, and which does not come into your mind daily without anyone's exhortation. But because, when I read your letters, I seem to hear you, and when I write to you I seem to speak with you, on that account I take greatest pleasure in your longest letter, and I myself in writing am often the longer.

\ciceroSection{46} This last thing I beg and exhort you: as good poets and industrious actors are wont to do, so be you most diligent in the last part and the conclusion of your duty and business, that this third year of your command may seem to have been, as the third act, the most perfect and the most adorned. You will most easily do this if you reckon that I --- whom you have always wished to please more, alone, than all men --- am always with you and present at all things which you say and do. It remains that I beg you, if you wish me and all yours to be well, that you most diligently serve your own health.
